부호에 따라 시프트 방향이 바뀌는 연산
\begin{equation}
f(v, s) = \begin{cases} v \ll s & (s > 0) \\ v \gg |s| & (s \le 0) \end{cases}
\label{eq:shift}
\end{equation}
은 고정소수점 신호처리나 양자화 신경망의 스케일링에서 자주 나타난다. 이 연산의 직관적인
구현은 부호 검사 분기를 포함하는데, 시프트 방향이 무작위에 가까우면 분기 예측이 불가능하므로
``분기를 제거하라(branchless)''는 것이 널리 알려진 최적화 상식이다\cite{hennessy,shiftop}.

그러나 이 상식은 분기 예측 실패가 비싸다는 전제 위에 서 있다. 사물인터넷 기기에 널리 쓰이는
Cortex-M0+\cite{m0plus}는 분기 예측기가 아예 없고 분기 비용이 1--3사이클로 고정된
프로세서다. 이런 코어에서도 분기 제거가 이득일까? 그리고 컴파일러에 최적화를 맡기는 것과
사람이 어셈블리를 직접 작성하는 것 사이에는 얼마나 차이가 있을까?

본 논문은 이 두 질문에 답하기 위해 RP2040 실제 하드웨어에서 세 가지 구현---분기 기준
구현(origin), 분기 제거 C 구현(shift\_pow2), 수동 어셈블리 구현(sum\_dual\_asm)---의 성능을
실측 비교한다. 실험 결과 분기 제거 C 구현은 오히려 기준보다 1.2배 느렸고, 동일한 분기 제거
아이디어를 어셈블리로 직접 작성하면 기준보다 1.5배 빨랐다. 우리는 컴파일러가 생성한 실제
코드를 분석하여 이 역전의 원인이 컴파일러의 능력 부족이 아니라 언어 규칙과 명령어 집합
구조가 만드는 구조적 한계임을 보이고, 이러한 한계를 대형 언어 모델(LLM)로 예측하고 보완하는
방향을 제시한다.
